% Two attachments and finite face error, 26 November 2026, input 5ddaabd76595.
% Continuation of 26 November 2026 at e2968b2898b5. Current proofs and primary
% passages are recorded in the prefix and curved-filling assessment in parent reviews.
% Prepared for Edward B. Baker III with OpenAI Codex assistance.
% Editorial execution, 26 November 2026. Actual serving model and effort
% unavailable. Research lead selection GPT-6.2 Astra/ultra with bounded
% GPT-6.2 Sol/ultra helper dispatches. Inputs at ece209cc09aa31629a8b496ae9d01957e71c70ac:
% EXACT_FACE_AND_BENDING.md dca779667aecc2fdf9b1565cbe8611462fb93c1b468f49a897eb3d642078b2dc
% APPROXIMATE_PHYSICAL_JETS.md f66c890a70717ad948356b9a268409765680abd7e7c2bf6e428bfda9ce53eab7
% COMPACT_MARKED_COLLAR_EXCLUSION.md ed02a13d0e724c5ea31fa6387c378c6a1ecf4ae25e78ae3250f1ed1594529422
% Scope: supplied-chart estimates and conditional applications to marked disks.
\documentclass[11pt,letterpaper]{article}
\usepackage[margin=1in]{geometry}
\usepackage[T1]{fontenc}
\usepackage{lmodern,microtype,amsmath,amssymb,amsthm,mathtools}
\usepackage{xurl}
\usepackage[hidelinks]{hyperref}
\setlength{\emergencystretch}{2em}
\numberwithin{equation}{section}
\newtheorem{theorem}{Theorem}[section]
\newtheorem{proposition}[theorem]{Proposition}
\newtheorem{lemma}[theorem]{Lemma}
\newtheorem{corollary}[theorem]{Corollary}
\DeclareMathOperator{\Cot}{Cot}
\DeclareMathOperator{\Ric}{Ric}
\DeclareMathOperator{\Rm}{Rm}
\DeclareMathOperator{\tr}{tr}
\newcommand{\HS}{\mathrm{HS}}
\newcommand{\rd}{\gamma}
\title{Deformation of marked spherical collars}
\author{Edward B. Baker III}
\date{Working manuscript, 26 November 2026}

\begin{document}
\maketitle
\begin{abstract}
A smooth conformally round umbilic three-sphere in a Riemannian
four-manifold admits a marked deformation to a round product on two
independently chosen short strips. In the normalized Gaussian chart
the signed cost of a fixed cutoff has a positive quadratic leading
term in each complete side and each chirality when the trace-free
second normal jet is nonzero. We estimate approximate face jets on
a supplied chart and obtain sharp inverse-width powers for this
interpolation on one nonround mode. For smoothly precompact marked
conformal pullbacks with positive definite limits, vanishing face
defect gives a subsequence with uniformly vanishing deformation
costs on both sides. On a hypothetical nonstandard homotopy sphere
with positive-class infimum below \(8\pi^2\), these comparisons
exclude an exact face at each sufficiently late charged disk and
exclude vanishing defect on a precompact subsequence of the marked
collars. A strictly convex Euclidean filling also treats certain nonumbilic
faces through their actual first jets and a positive exterior quadratic form.
\end{abstract}

\section{Physical collars and chiral costs}

All manifolds, metrics, charts and embeddings are smooth. Metrics
are Riemannian and the ambient dimension is four. On an oriented
source \(M\), let \(W^\sigma_g\) be the Weyl curvature operator on
the real bundle \(\Lambda^\sigma\) of two-forms, for
\(\sigma\in\{+,-\}\). Use the curvature convention for which the unit
sphere has the identity curvature operator. Set
\[
 E_\sigma(g;U)=\int_U |W^\sigma_g|_{\HS}^2\,dV_g .
\]
The four-index tensor norm satisfies
\(\lvert W\rvert_{\mathrm{tens}}^2
 =4(\lvert W^+\rvert_{\HS}^2+\lvert W^-\rvert_{\HS}^2)\).
In dimension four \(g\mapsto e^{2u}g\) fixes the Hodge star on
two-forms and scales the Weyl operator by \(e^{-2u}\) and volume
by \(e^{4u}\). Both chiral measures are preserved on every physical
subset. On a closed homotopy four-sphere the joined energies agree
by signature. Regional energies are measured separately in the
ambient orientation \cite{Hirzebruch1971Signature,Baker2026ChiralMaster}.

Take a compact source with a separating embedded cooriented
three-sphere \(\Sigma\), its specified angular marking
\(a:S^3\to\Sigma\) and a physical bicollar \(\Phi\) extending \(a\).
Write \(P_L,P_R\) for the two complete closed sides and let \(\nu\)
point from \(P_L\) to \(P_R\). At the cut their outward normals
are \(\nu\) and \(-\nu\), respectively. A path \(g_s\) starting
at \(g\) gives signed changes
\[
 \Delta_{j,\sigma}(s)=E_\sigma(g_s;P_j)-E_\sigma(g;P_j),
 \qquad j=L,R .
\]
The joined change is their sum. For a cutoff supported in one
strip on each side, the signed changes are exactly the differences
of deformed and original strip energies. The physical marking and
its inverse determine where these strips and the capped pieces
attach. A change in marking must be carried through residual
attaching tubes and the gluing maps, as in the companion analysis
of marked handle cuts \cite{Baker2026HandleCuts,Baker2026CollarProgram}.

\section{Exact first-order matching}\label{sec:exact}

Let \(\rd\) be the unit round metric on the marked \(S^3\). On
\(\Sigma\) write \(h=g|_{T\Sigma}\),
\(K(X,Y)=g(\nabla_X\nu,Y)\), and \(H=\tr_h K\).
Suppose \(a^*h=e^{-2v}\rd\) for some smooth \(v\) in the given
marking and \(K=Hh/3\). In an original signed \(g\)-Gaussian chart
\(\Psi(x,q)=\exp^g_{a(x)}(q\nu)\), choose
\[
 u(\Psi(x,q))=\rho(q)\bigl(v(x)-qH(a(x))/3\bigr),
\]
where \(\rho\) is one near zero and supported inside a fixed compact
subchart. Extend \(u\) by zero. For \(\widehat g=e^{2u}g\) the
conformal connection formula gives
\[
 \widehat\nu=e^{-u}\nu,\qquad
 \widehat K=e^u\bigl(K+\nu(u)h\bigr).
\]
In fact the tangent part of \(\widehat\nabla_X(e^{-u}\nu)\) is
\(e^{-u}(\nabla_X\nu+\nu(u)X)\) for tangent \(X\). The chosen value
of \(u\) gives the round induced metric, and its independently chosen
normal derivative makes the face totally geodesic. The average mean
curvature convention in \cite[\S6.1, Proposition 6.1 and Remark 6.5]{CurryGover2015Introduction}
corresponds to \(H/3\) here.

Take the signed normal exponential chart for \(\widehat g\),
\[
 F(x,t)=\exp^{\widehat g}_{a(x)}(t\widehat\nu),\qquad
 F^{-1}(p)=\bigl(a^{-1}(\pi_{\widehat g}(p)),t_{\widehat g}(p)\bigr),
\]
where \(\pi_{\widehat g}\) is the normal foot and \(t_{\widehat g}\)
is signed distance on the tube. In the original Gaussian chart,
\(q(F(x,t))=e^{-v(x)}t+O(t^2)\) uniformly on the compact face.
For any other signed collar parameter \(q\) the leading factor is
\(dq(\widehat\nu)>0\) there. Thus both normalized widths can be
chosen within the physical collar, separately on each side.
In the overlap the transition maps are \(\Phi^{-1}F\) and
\(F^{-1}\Phi\). At a normalized height \(t\) the marking is
\(F(\,\cdot\,,t)\), with inverse the first component of \(F^{-1}\).
The supplied outer-face markings are retained. The two copies
\(a_L,a_R\) of \(a\) give the reassembly \(a_Ra_L^{-1}\)
and its inverse \(a_La_R^{-1}\), reversing the boundary
orientations.

Write \(F^*\widehat g=dt^2+h_t\), so \(h_0=\rd\) and
\(\partial_t h_0=0\). Fix smooth profiles \(0\le\chi_j\le1\)
on \([0,1]\), one near zero and zero near one. Write \(t=\epsilon_j w_j r\)
on each half, where \(\epsilon_L=-1\), \(\epsilon_R=1\),
\(0\le r\le1\), and \(0<w_j\le1\). Set
\begin{equation}
 D_t=h_t-\rd,\qquad h_{t,s}=\rd+(1-s\chi_j(r))D_t,
 \qquad 0\le s\le1 .
 \label{eq:exact-path}
\end{equation}
Both profiles are constant and agree near the middle, so the
interpolation is smooth there. Convexity gives positivity. Near
each outer face the interpolant equals the normalized source.
Let \(N_s\) be the pushforward, extended by \(\widehat g\). The
physical path \(g_s=e^{-2u}N_s\) starts at \(g\), agrees with it
off the two strips and ends in a product germ up to conformal
scale. The path \(e^{2s\zeta u}g_s\), where \(\zeta=1\) near the
middle and is supported inside the terminal product germ, exposes
an exact product at its physical endpoint. The chiral costs of
the two paths agree on each full side.

\begin{theorem}[Exact-face cost and peak]\label{thm:exact}
For the fixed source, both complete sides and both signs satisfy
\(\sup_s|\Delta_{j,\sigma}(s)|\le C_j w_j\).
Let \(A=\tfrac12\partial_t^2h_t|_{t=0}\) and take its trace-free
part \(A^\circ\) in \(\rd\). Then
\begin{equation}
 \Delta_{j,\sigma}(s)=w_j s^2 c_j+O_{C_s^2}(w_j^2),\qquad
 c_j=\frac{\|A^\circ\|_{L^2(\rd)}^2}{16}
       \int_0^1 r^4(\chi_j''(r))^2\,dr .
 \label{eq:exact-coefficient}
\end{equation}
The remainder is uniform through two \(s\)-derivatives. If
\(A^\circ\not\equiv0\), then \(c_j>0\) on both sides. For
sufficiently small independent widths each side and each joined
energy peaks at \(s=1\), the latter with value
\(E_\sigma(g;M)+w_Lc_L+w_Rc_R+O(w_L^2+w_R^2)\).
\end{theorem}

\begin{proof}
Taylor expansion gives
\(\|(\nabla^\rd)^k\partial_t^p D_t\|\le C|t|^{2-p}\) for
\(p+k\le2\). Write \(b=1-s\chi_j(r)\). Its derivatives are
\(b'=-s\epsilon_j\chi_j'/w_j\) and \(b''=-s\chi_j''/w_j^2\).
The first and second normal derivatives of \(bD_t\) are
\(bD_t'+b'D_t\) and \(bD_t''+2b'D_t'+b''D_t\).
Angular derivatives commute with \(b\). These bounds, together
with the fixed ellipticity interval, give uniformly bounded
curvature and density. The energy in either strip, throughout the
path and at the source, is \(O(w_j)\).

For the coefficient use \(D_t=t^2A+O(t^3)\) with differentiated
Taylor bounds. At \(t=\epsilon_jw_jr\) this gives
\[
 \tfrac12\partial_t^2h_{t,s}=(1-sB_j(r))A+O(w_j),
 \qquad B_j=\tfrac12(r^2\chi_j)''.
\]
For the metric \(dt^2+h_{t,s}\) the normal curvature block is
\(-\partial_t K+Kh^{-1}K\) with \(K=\partial_t h_{t,s}/2\).
Since \(K=O(w_j)\), its leading term is \(-(1-sB_j)A\).
Codazzi gives a mixed block \(O(w_j)\). Gauss gives a tangential
block equal to the round identity plus \(O(w_j^2)\), using
\(h_{t,s}=\rd+O(w_j^2)\) through two angular derivatives and
\(K=O(w_j)\). Passage to metric-orthonormal angular frames
uses \(I+O(w_j^2)\). For \(\alpha_i=dt\wedge e^i\) in such
a frame, \((\alpha_i+\sigma*\alpha_i)/\sqrt2\) gives the
chiral basis. Its curvature block is the half-sum of the
normal and tangential blocks plus signed mixed terms. Taking
its trace-free part gives
\[
 W^\sigma=-\tfrac12(1-sB_j)A^\circ+O(w_j),\qquad
 dV_{N_s}=(1+O(w_j^2))\,|dt|\,dA_\rd .
\]
The two-jet is affine in \(s\), and all operations forming these
blocks and densities are smooth on the fixed ellipticity interval.
Thus the errors have the same estimates through the first two
\(s\)-derivatives.

Abbreviate \(\chi=\chi_j\), \(B=B_j\). Since
\((r^2\chi)'=0\) at both ends, \(\int_0^1B\,dr=0\).
In the square of \(B=r^2\chi''/2+2r\chi'+\chi\), the cross
terms satisfy
\[
 \begin{aligned}
 \int_0^1 2r^3\chi'\chi''\,dr&=-3\int_0^1r^2(\chi')^2\,dr,\\
 \int_0^1 r^2\chi\chi''\,dr
     &=\int_0^1\chi^2\,dr-\int_0^1r^2(\chi')^2\,dr,\\
 \int_0^1 4r\chi\chi'\,dr&=-2\int_0^1\chi^2\,dr .
 \end{aligned}
\]
All boundary terms vanish by the cutoff conditions and the
powers of \(r\). Hence \(\int B^2=\frac14\int r^4(\chi'')^2\),
which is positive since \(B=1\) near zero. The difference of
the leading squared Weyl norms integrates to
\[
 \tfrac14w_j\|A^\circ\|_2^2
       \int_0^1\bigl((1-sB_j)^2-1\bigr)\,dr .
\]
The linear term vanishes and the quadratic term gives
\eqref{eq:exact-coefficient}. If \(c_j>0\), then
\(\partial_s^2\Delta_{j,\sigma}=2w_jc_j+O(w_j^2)>0\)
and \(\Delta_{j,\sigma}(1)>\Delta_{j,\sigma}(0)=0\).
Convexity proves the endpoint maximum for each side and the sum.
\end{proof}

If the jet is written as \(J^\circ=(\partial_t^2h_t|_{0})^\circ\),
the factor \(1/16\) in \eqref{eq:exact-coefficient} becomes \(1/64\).
The reverse path from the product has signed change
\(\Delta_{j,\sigma}(1-s)-\Delta_{j,\sigma}(1)\), with leading
term \(w_jc_j((1-s)^2-1)\). Reversing the normal exchanges the
side labels, widths and profiles. Reversing the ambient orientation
exchanges the chiralities. For \(A^\circ=0\), the \(O(w_j^2)\)
bound holds without selection of a maximizer. The constants here
are for a fixed source, normalized chart and profiles.

\section{Approximate jets in a supplied Gaussian chart}

Now take an actual marked physical chart with metric
\(dt^2+h_t\) on two strips of independent widths \(0<w_j\le1\).
Use reference \(\bar p=dt^2+\rd\), \(u_t=h_t-\rd\), and
\(p_s=\bar p+(1-s\chi_j(r))u_t\). Assume
\(\lambda\rd\le h_t\le\Lambda\rd\) on both strips. The angular
norm \(\|\cdot\|_{C^2_\rd}\) sums supremum norms for the value
and first two angular derivatives. Put
\[
 \begin{gathered}
 \alpha=\|u_0\|_{C^2_\rd},\qquad
 \beta=\|\partial_tu_0\|_{C^2_\rd},\qquad
 M_j=\sup_{t\ \mathrm{on\ side}\ j}\|\partial_t^2h_t\|_{C^2_\rd},\\
 Z_j=\frac{\alpha}{w_j^{3/2}}+\frac{\beta}{\sqrt{w_j}}+\sqrt{w_j}M_j,
 \qquad A_\rd=2\pi^2 .
 \end{gathered}
\]
The path is a smooth intrinsic deformation that agrees with the
source near both outer faces and ends in an exact product germ.

\begin{theorem}[Two physical widths]\label{thm:approx}
On each half-strip \(U_j\), for both signs and all \(0\le s\le1\),
\begin{equation}
 E_\sigma(p_s;U_j)^{1/2}
 \le C\sqrt{A_\rd}\bigl(Z_j+w_j^{3/2}Z_j^2\bigr).
 \label{eq:approx}
\end{equation}
The constant depends on dimension, reference curvature, the bounds
\(\min(1,\lambda)\), \(\max(1,\Lambda)\), and the fixed cutoff
derivative bounds. It is independent of the ratio of widths.
If \(C_j^{\rm app}\) is the square of the right side, the signed
increment lies in
\([-E_\sigma(g;U_j),\,C_j^{\rm app}-E_\sigma(g;U_j)]\).
The joined peak is at most \(E_\sigma(g;M)+C_L^{\rm app}+C_R^{\rm app}\).
\end{theorem}

\begin{proof}
Write \(w=w_j\), \(v=(1-s\chi_j)u_t\),
\(L_k=\|\chi_j^{(k)}\|_\infty\),
\(P=\alpha+w\beta+w^2 M_j/2\), and \(Q=\beta+w M_j\).
The latter two bound \(u_t,\partial_tu_t\) in the angular norm.
For \(D=\nabla^\rd\) the derivatives through total order two are
\[
 \begin{aligned}
 D^k v&=(1-s\chi_j)D^ku_t &&(0\le k\le2),\\
 \partial_t D^k v&=(1-s\chi_j)\partial_tD^ku_t
       -s\epsilon_j\chi_j'D^ku_t/w &&(0\le k\le1),\\
 \partial_t^2 v&=(1-s\chi_j)\partial_t^2u_t
       -2s\epsilon_j\chi_j'\partial_tu_t/w-s\chi_j''u_t/w^2 .
 \end{aligned}
\]
Thus the zero-, one-, and two-normal bounds are \(P\),
\(Q+sL_1P/w\), and \(M_j+2sL_1Q/w+sL_2P/w^2\) in the
indicated angular norms. The full product jets through order two
are bounded, by order, by \(Cw^{3/2}Z_j\), \(Cw^{1/2}Z_j\),
and \(Cw^{-1/2}Z_j\).

In oriented \(\bar p\)-orthonormal frames identify the metric
orthonormal frames via the positive inverse square root.
Let \(F_\sigma(v)\) be the chiral Weyl matrix for \(\bar p+v\)
in that frame multiplied by \((dV_{\bar p+v}/dV_{\bar p})^{1/2}\).
Its squared \(L^2(dV_{\bar p})\) norm gives the energy. This is
the density convention of the first-order splice in
\cite{Baker2026ChiralMaster}. In the product connection
\(\overline\nabla\), the connection difference is a contraction
of \((\bar p+v)^{-1}\overline\nabla v\). One derivative and the
square of this difference give curvature terms linear in
\(\overline\nabla^2v\) and quadratic in \(\overline\nabla v\).
The projection and density are smooth functions of the metric.
Since \(\bar p\) is conformally flat, the linearization and remainder
obey
\[
 \begin{aligned}
 |\mathcal L_\sigma v|&\le
     C(|\overline{\Rm}|\,|v|+|\overline\nabla^2v|),\\
 |F_\sigma(v)-\mathcal L_\sigma v|&\le
     C\bigl(|\overline{\Rm}|\,|v|^2+
            |v|\,|\overline\nabla^2v|+|\overline\nabla v|^2\bigr).
 \end{aligned}
\]
The terms with \(\overline{\Rm}\) come from zeroth-order
coefficient changes. The quadratic first-jet contribution includes
the square of \(sL_1P/w\) and its cross terms. Using ellipticity
to bound \(|v|\) in the mixed quadratic term yields the pointwise
estimate \(C(w^{-1/2}Z_j+w Z_j^2)\). The reference volume
\(A_\rd w\) gives \eqref{eq:approx}. The full-side changes
subtract the raw strip energies since the metric is fixed
off the two strips.
\end{proof}

Uniform ellipticity and \(Z_L,Z_R\to0\) imply vanishing allowances.
For shrinking widths it is enough to have, for
\(w_{\min}=\min(w_L,w_R)\),
\[
 \alpha=o(w_{\min}^{3/2}),\qquad \beta=o(w_{\min}^{1/2}),\qquad
 \sqrt{w_j}M_j\longrightarrow0\quad(j=L,R).
\]
For exact face jets with bounded \(M_j\), the squared bound gives
\(O(w_j)\). The hypotheses are on physical Taylor data, and any
preliminary metric adjustment outside the conformal class or
added connector must be included in the joined and side budgets.

\section{Sharp powers for a fixed tangential mode}

Identify the marked round \(S^3\) with \(SU(2)\), and let
\((\omega^1,\omega^2,\omega^3)\) be a unit round left-invariant
coframe, with cyclic brackets \([e_j,e_k]=2e_i\) for the dual
frame. Use the fixed tensor
\[
 Q=\omega^1\otimes\omega^1-\omega^2\otimes\omega^2,\qquad
 h_t=\rd+(a+bt)Q .
\]
Choose amplitudes to make the metric uniformly positive on the
given chart. The tensor has eigenvalues \(q_i=(1,-1,0)_i\), is
trace-free and satisfies \(\|Q\|_2^2=4\pi^2\). To check the
nonroundness, for \(h=\rd+fQ\) with constant \(f\) put
\(l_i=1+fq_i\) and \(\mu_i=2l_i/\sqrt{l_1l_2l_3}\).
In the orthonormal frame \(E_i=l_i^{-1/2}e_i\), the cyclic
bracket coefficients are \(\mu_i\). The Koszul formula gives
\(2\langle\nabla_{E_i}E_j,E_k\rangle=\mu_k-\mu_i+\mu_j\)
for a cyclic triple. Substitution into the curvature formula
gives the Ricci eigenvalues
\(\tfrac12(\mu_i^2-(\mu_j-\mu_k)^2)=2+4f q_i+O(f^2)\)
for \(\{i,j,k\}=\{1,2,3\}\). They differ for sufficiently
small nonzero \(f\).

On the strip of width \(w=w_j\), the interpolant has tangential
coefficient \(f_s=(1-s\chi_j(r))(a+bt)\). Its linearized
two-normal curvature is \(-f_s''Q/2\), with primes here taken
in \(t\). The tangential and mixed blocks carry no two-normal
linear terms. The trace-free chiral half-sum used in
Theorem~\ref{thm:exact} therefore gives
\(-f_s''Q/4\) as the principal term of \(F_\sigma(f_s Q)\)
in either sign. The density normalization leaves this linear
term unchanged. Differentiation at the native width gives
\[
 f_s''=-s\left[
  \frac{a}{w^2}\chi_j''(r)
  +\frac{\epsilon_j b}{w}\bigl(r\chi_j''(r)+2\chi_j'(r)\bigr)
 \right].
\]
Define
\begin{equation}
 J_j(a,b)=\frac{\|Q\|_2^2}{16}\int_0^1
 \left[
  \frac{a}{w_j^{3/2}}\chi_j''
  +\frac{\epsilon_j b}{\sqrt{w_j}}\bigl(r\chi_j''+2\chi_j'\bigr)
 \right]^2\,dr .
 \label{eq:gram}
\end{equation}
The profiles \(\chi_j''\) and \((r\chi_j)''\) have a positive
definite Gram matrix. Otherwise \(((A+Br)\chi_j)''=0\) for
a nonzero \((A,B)\). The affine function \((A+Br)\chi_j\)
vanishes near one and equals \(A+Br\) near zero, a contradiction.
Consequently \(J_j\) is comparable to \(a^2/w_j^3+b^2/w_j\),
with constants for the fixed profile and \(Q\). The mixed term
carries the factor \(\epsilon_j ab\). Its profile coefficient
is positive since
\(\int_0^1\chi_j''(r\chi_j''+2\chi_j')\,dr
=\int_0^1r(\chi_j'')^2\,dr\), the boundary term
\([(\chi_j')^2]_0^1\) vanishing.

\begin{proposition}[Signed cost of the mode]\label{prop:mode}
For \(\delta_j=|a|+w_j|b|\) sufficiently small, the complete
side changes in either sign satisfy
\begin{equation}
 \Delta_{j,\sigma}(s)=s^2 J_j(a,b)
             +O\bigl((w_j+\delta_j)J_j(a,b)\bigr) ,
 \label{eq:mode-sign}
\end{equation}
uniformly in \(s\) and the ratio of widths. The leading joined
peak increment is \(J_L+J_R\), with additive error bounded by
the sum of the two errors in \eqref{eq:mode-sign}.
For \(a=0\) and fixed \(b\), uniformly for \(0\le s\le1\),
\[
 \lim_{w_j\to0} w_j\Delta_{j,\sigma}(s)=
 \frac{s^2 b^2\|Q\|_2^2}{16}
 \int_0^1 (r\chi_j''+2\chi_j')^2\,dr .
\]
\end{proposition}

\begin{proof}
The fixed angular tensor and the density expansion in the proof
of Theorem~\ref{thm:approx} bound the error after removal of
\(-f_s''Q/4\) pointwise by
\[
 \begin{aligned}
 C\bigl(|f_s|+|f_s'|+|f_s|\,|f_s''|+|f_s'|^2\bigr)
 &\le C(w_j+\delta_j)\delta_j/w_j^2 .
 \end{aligned}
\]
Indeed the three jets in \(t\) are bounded in order by
\(C\delta_j,C\delta_j/w_j,C\delta_j/w_j^2\).
The mixed and tangential linear terms use at most one normal
derivative. The squared reference norm of the principal part is
\(s^2J_j\), and the error norm is bounded by
\(C(w_j+\delta_j)\sqrt{J_j}\). For the incoming metric, its
coefficient \(a+bt\) has second normal derivative zero.
The same expansion bounds its energy on the strip by
\[
 Cw_j(\delta_j+|b|+b^2)^2
 \le C(w_j+\delta_j)^2J_j .
\]
The \(b^2\) term controls the quadratic first jet also for large
\(|b|\) on shrinking charts with \(\delta_j\) small. Squaring the
outgoing density norm and subtracting the incoming energy proves
\eqref{eq:mode-sign}. The uniform bound and the value at \(s=1\)
give the leading peak, and the pure first-jet limit follows by
substituting \(a=0\) and \(\delta_j=w_j|b|\).
\end{proof}

For fixed small \(a\ne0\) and fixed \(b\), an exact
\(w_j^{-3}\) coefficient retains the quadratic normal-curvature
terms. For a given profile \(\chi_j\), put
\[
 \begin{aligned}
 l_i(r)&=1+a(1-s\chi_j(r))q_i, & k_i&=\frac{l_i'}{2l_i},\\
 N_i&=-\frac{l_i''}{2l_i}+\frac{(l_i')^2}{4l_i^2}, &
 T_i&=-k_m k_n \quad(\{i,m,n\}=\{1,2,3\}),
 \end{aligned}
\]
now with primes in \(r\).
Before taking the limit the same diagonal lengths have the
extra term \((1-s\chi_j)\epsilon_j b w_j r q_i\), which
converges to zero through two \(r\)-derivatives uniformly in \(s\).
For a tangential diagonal length \(\widetilde l_i(t)\), the
normal block in a metric orthonormal frame has diagonal
\(-\widetilde l_{i,tt}/(2\widetilde l_i)
 +(\widetilde l_{i,t})^2/(4\widetilde l_i^2)\).
The tangential block has the intrinsic sphere curvature and
the Gauss term
\(-(\widetilde l_{m,t}/(2\widetilde l_m))
  (\widetilde l_{n,t}/(2\widetilde l_n))\).
After multiplying by \(w_j^2\), these converge to \(N_i\) and
\(T_i\), respectively. Intrinsic curvature is bounded, and the
mixed curvature is \(O(w_j^{-1})\) by Codazzi, so neither
contributes to the rescaled limit. The chiral diagonal
entries converge to \((N_i+T_i)/2\), before subtracting their
mean. With the volume factor this gives
\begin{equation}
 \begin{aligned}
 w_j^3\Delta_{j,\sigma}(s)&\longrightarrow\Lambda_j(s,a),\\
 \Lambda_j(s,a)
 &=\frac{A_\rd}{4}\int_0^1\sqrt{l_1l_2l_3}
   \sum_{i=1}^3
   \left(N_i+T_i-\frac13\sum_{m=1}^3(N_m+T_m)\right)^2\,dr.
 \end{aligned}
 \label{eq:fixed-a}
\end{equation}
The incoming energies vanish at the indicated rescaling.
Convergence is uniform in \(s\) and identical in the two signs.
If the same profile is used in both strips, the functions
\(\Lambda_j\) agree. The sign \(\epsilon_j\) in a first
normal derivative cancels in these quadratic terms.

The formulas for \(l_i'\) and \(l_i''\) exhibit a factor \(sa\)
in each leading curvature entry. At \(a=0\) after division by that
factor, the trace-free part of \(N_i+T_i\) has entries
\(\chi_j''q_i/2\). Thus
\(\Lambda_j(s,a)=s^2a^2 G_j(s,a)\), where \(G_j\) is smooth
through \(s=0,a=0\) and
\[
 G_j(s,a)=\frac{\|Q\|_2^2}{16}\int_0^1(\chi_j'')^2\,dr+O(|a|),
 \qquad \partial_s G_j=O(|a|).
\]
The constant term is positive. For sufficiently small nonzero
\(a\), differentiation shows that \(\Lambda_j(s,a)\)
increases for \(s>0\). The leading joined peak for independently
shrinking widths is therefore
\[
 w_L^{-3}\Lambda_L(1,a)+w_R^{-3}\Lambda_R(1,a)
      +o(w_L^{-3}+w_R^{-3}) .
\]
In particular the powers \(w^{-3}\) and \(w^{-1}\) both occur
for the specified interpolation. These are sharpness tests on
a single nonround tensor mode in a supplied chart. They do
not establish lower bounds for other deformations or charts.

\section{Small face defect in a compact family}

On a smooth cooriented three-sphere with first and second
fundamental forms \(h,K\), define
\[
 \mathcal D_g(\Sigma)=
 \int_\Sigma\bigl(|\Cot(h)|_h+|K^\circ|_h^3\bigr)\,dA_h .
\]
Here \(\Cot\) is the three-covariant intrinsic Cotton tensor.
In dimension three it is conformally invariant and vanishes
precisely for locally conformally flat metrics. Its norm has
weight \(-3\). The norm of \(K^\circ\) has weight \(-1\) under
ambient conformal change, while \(dA_h\) has weight \(3\).
Thus both integrals are conformally invariant
\cite[\S4.1 and \S6.1]{CurryGover2015Introduction}.

\begin{theorem}[Smooth marked precompactness]\label{thm:compact}
Let \(\Phi_i:S^3\times[-1,1]\to M_i\) be smooth embedded
bicollars of specified cooriented separating spheres \(\Sigma_i\),
with physical metrics \(g_i\) and fixed markings of their complete
sides. Suppose positive conformal multiples of \(\Phi_i^*g_i\)
are precompact in the smooth topology of this common closed slab,
with positive definite limits. If \(\mathcal D_{g_i}(\Sigma_i)\to0\),
a subsequence admits marked paths to round product germs up to
conformal scale, with \(o(1)\) absolute signed changes in both
chiralities on both complete sides, uniformly along the paths.
If an application requires separate protected regions to remain
fixed, the strips must fit inside the permitted region with the
width bounds in the proof.
\end{theorem}

\begin{proof}
Pass to \(p_i=e^{2u_i}\Phi_i^*g_i\to p_\infty\) smoothly.
The two nonnegative face integrands converge on the compact
sphere. Vanishing of their limiting integrals gives an umbilic
face and \(\Cot(h_\infty)=0\). By conformal development the
smooth compact simply connected locally conformally flat
sphere is conformally round \cite{Kuiper1949}. Its developing
map is a local diffeomorphism whose image is open and closed
in \(S^3\). Compactness makes it a covering, which is a
diffeomorphism since \(S^3\) is simply connected. Write the
round representative as \(\eta=e^{2v_0}h_\infty=\vartheta^*\rd\)
on the original sphere, choosing \(\vartheta\) orientation
preserving by a round reflection if necessary.

Use the normalization of Section~\ref{sec:exact} for the limit, with face
value \(v_0\) and normal derivative \(-H_\infty/3\) of a
fixed function \(v\). Its signed Gaussian chart \(F\) on a
smaller closed slab fixes the original angular coordinate at
the center. Put \(b=dt^2+\eta\) and
\(m_\infty=F^*(e^{2v}p_\infty)=dt^2+h_t\) there. The limiting
face has \(h_0=\eta\) and \(\partial_t h_0=0\). The actual
\(m_i=F^*(e^{2v}p_i)\) need not be Gaussian for this same
\(F\). With the full product covariant norm let
\(\delta_i=\|m_i-m_\infty\|_{C^2(b)}\). Choose subwidths
\(w_{i,L}=(\delta_i+i^{-1})^{1/3}\) and \(w_{i,R}=w_{i,L}/2\)
for late \(i\). More generally two available widths tending
to zero suffice if \(\delta_i/w_{i,j}^2\) stays bounded.
The displayed choice satisfies this condition and fits inside
the smaller slab for late indices. A fixed positive clearance
in the permitted region suffices if there are extra protected sets.

On each side use the interpolation of the full four-tensor
\[
 m_{i,s}=b+(1-s\chi_j(|t|/w_{i,j}))(m_i-b).
\]
The first three orders of the limiting mismatch have
bounds \(O(w^2),O(w),O(1)\) for width \(w=w_{i,j}\).
For \(m_i-m_\infty\) the corresponding cutoff jet bounds
are \(C\delta_i\), \(C(1+w^{-1})\delta_i\) and
\(C(1+w^{-1}+w^{-2})\delta_i\). This gives bounded jets
through order two for all normal, mixed and tangential
components. The path has uniform ellipticity by convexity.
Thus both strip energies, including their incoming values,
are \(O(w_{i,j})\) in either sign, uniformly along the path.

Push forward by \(\Phi_i F\) with inverse \(F^{-1}\Phi_i^{-1}\)
on its image and restore the inverse conformal multiplier.
This returns the source metric near both outer faces, with zero
gauge charge. If needed, an endpoint conformal change supported
in a smaller product germ exposes its product representative.
For \(\alpha_i(x)=\Phi_i(x,0)\) the round cap coordinate
\(z=\vartheta(x)\) marks the actual face by
\(z\mapsto\alpha_i(\vartheta^{-1}(z))\), with inverse
\(\vartheta\alpha_i^{-1}\). This records the angular change in
both directions.
\end{proof}

For a parametrized disk \(A_i:B^4\to D_i\) the transition from
the cap boundary to the disk coordinates is
\((A_i|_{S^3})^{-1}\alpha_i\vartheta^{-1}\), with inverse
\(\vartheta\alpha_i^{-1}(A_i|_{S^3})\).
Every sphere diffeomorphism is isotopic to an orthogonal map of
its orientation sign \cite[p.~553]{Hatcher1983Smale}.
Given an isotopy \(G_s\) from that linear boundary map to the
transition, extend by \((r,z)\mapsto(r,G_{\rho(r)}(z))\)
with \(\rho=0\) near zero and \(\rho=1\) near one.
The inverse uses \(G_{\rho(r)}^{-1}\). Near the center the
extension is linear. Choosing the cap orientation so the
boundary gluing reverses the induced orientations gives smooth
oriented reassembly relative to the retained exterior away
from the joining collar. A trace application must instead
transport all its marked attachments under the angular change
\cite{Baker2026HandleCuts}.

\section{The joined threshold and charged disks}

Let \(M\) be a closed oriented smooth homotopy four-sphere.
Write \(T=4\pi^2\) and
\[
 I=I^+(M)=\inf\{E_+(h;M):Y([h])>0\}.
\]
The empty infimum is \(+\infty\). For the positive metrics
considered below the infimum is finite, and either chirality gives
the same closed energy by signature. We use the master's positive
recognition criterion \(E_+\le T\) and its charged-disk selection
on a positive subcritical minimizing sequence
\cite[Recognition and charged-disk theorems]{Baker2026ChiralMaster}.

\begin{lemma}[Strict joined bound and positive caps]\label{lem:cap}
Suppose \(Y([g])>0\) and a smooth joined metric path from \(g\)
has peak \(H_0=\sup_s E_+(g_s;M)<2T\). Suppose its endpoint
\(h\) has a marked round product germ up to smooth positive
conformal factor at a separating sphere. Then either complete
side can be capped by a smooth marked ball, with positive Yamabe
class and exactly its endpoint side energy in each chirality.
\end{lemma}

\begin{proof}
If the first conformal-Laplacian eigenvalue becomes zero, its
positive eigenfunction gives a scalar-flat metric \(q\) in that
class. In the present normalization, Gauss--Bonnet and signature
give \(E_+(q)=8\pi^2+\frac14\int_M|\Ric_q|^2\,dV_q\).
The first eigenvalue is continuous along the metric path and its
sign is the Yamabe sign. Consequently the path stays positive.
Gauss--Bonnet on a positive volume-one Yamabe representative
also gives \(Y\ge\sqrt{48(2T-H_0)}\) along the path
\cite[\S2]{LeeParker1987Yamabe}
\cite[A short marked cut]{Baker2026ChiralMaster}.

Expose the product near the cut by a supported conformal change
at no cost. In a smaller exact product replace the normal
length element \(dt\) by \((1+L\varphi(t))\,dt\), with \(L\ge0\)
and a nonnegative interior bump \(\varphi\) of integral one.
Arclength identifies the changed region with a longer product.
The metrics on the joined source are joined to \(h\) by the same
stretching path and the added length has zero chiral energy.
The positive bound on \(Y\) therefore holds uniformly in \(L\).
Every compactly supported test on a retained side with a
half-infinite product end extends to a sufficiently stretched
joined metric with the same Yamabe quotient. In outward product
arclength \(\ell\) a positive conformal scale equal to \(e^{-\ell}\)
for large \(\ell\) compactifies the end in the angular marking
using \(r=e^{-\ell}\), with inverse \(\ell=-\log r\) away from
the center. The cap is flat near the added point and conformally
flat throughout. Choose the scale to be one near the joining
sphere. Smooth tests avoiding the point are dense in \(H^1\), since
a radial cutoff of a fixed smooth test has squared gradient
cost \(O(r^2)\) near the added point. Conformal covariance gives
the same strictly positive lower bound for the capped Yamabe
class. There is zero cap energy in either sign.
\end{proof}

Fix a smooth positive source \(g_i\) with joined energy
\(e_i<2T\), and let \(d_i=e_i-I\). Let \(D_i\) be an
actual parametrized standard closed disk with
\(V_i=M\setminus\operatorname{int}D_i\). For either sign set
\(m_i^\sigma=E_\sigma(g_i;D_i)\) and
\(v_i^\sigma=E_\sigma(g_i;V_i)\), so \(m_i^\sigma+v_i^\sigma=e_i\).
For a path as in Lemma~\ref{lem:cap} at its boundary, write
\(\Delta_{V,i}^\sigma=E_\sigma(h_i;V_i)-v_i^\sigma\) for its
signed complete-exterior endpoint change. Its exterior cap,
identified with \(M\) by the disk parametrization and the
sphere transition extension, is an actual positive metric on
the same oriented smooth type. Consequently
\begin{equation}
 \begin{aligned}
 E_\sigma(\widehat g_{V,i};M)
   &=e_i-m_i^\sigma+\Delta_{V,i}^\sigma\ge I,\\
 m_i^\sigma&\le d_i+\Delta_{V,i}^\sigma,\qquad
 \Delta_{V,i}^+-\Delta_{V,i}^-=m_i^+-m_i^- .
 \end{aligned}
 \label{eq:disk-budget}
\end{equation}
The last identity uses equality of the two closed capped energies.
If \(G_i^\sigma\) is an upper allowance on the signed exterior
change, then \(m_i^\sigma>d_i+G_i^\sigma\) contradicts existence
of the path at that finite index. If
\(e_i-m_i^\sigma+G_i^\sigma\le T\), the capped metric gives
direct recognition. In these comparisons the peak consumes both
strips and any connectors, and the capped exterior consumes its
entire changed side \cite[Exterior replacement and marked cuts]{Baker2026ChiralMaster}.

\begin{corollary}[A fixed exact face on a physical disk]\label{cor:fixed}
For a fixed positive source as above, suppose the disk boundary
is conformally round in some smooth angular marking and umbilic,
and admits arbitrarily narrow strips on both sides in the permitted
physical bicollar. Then \(m_i^+=m_i^-\le d_i\).
\end{corollary}

\begin{proof}
Theorem~\ref{thm:exact} gives joined peak at most
\(e_i+C_i(w_D+w_V)\) and \(|\Delta_{V,i}^\sigma|\le C_iw_V\)
with a fixed constant for this metric and chart. At this \(i\)
choose both positive widths narrow enough to stay strictly below
\(2T\). Perform the cap construction for each sufficiently small
choice, then send the widths to zero while holding \(g_i,D_i\),
the conformal normalization, the chart and \(C_i\) fixed.
Both inequalities and the identity in \eqref{eq:disk-budget}
give the conclusion. A conformally round marking different
from the disk's original boundary coordinates is admissible by
the sphere-extension argument following Theorem~\ref{thm:compact}.
\end{proof}

The order of quantifiers is material. For a hypothetical
nonstandard \(M\) with \(I<2T\), the master supplies a fixed
sign \(\sigma_*\), a subsequence of unit-volume positive Yamabe
representatives \(g_i\) of a minimizing sequence with \(e_i\to I\),
original physical marked disks
\(D_i\), and constants \(a,b>0\) with
\[
 m_i^{\sigma_*}\ge a,\qquad v_i^{\sigma_*}\ge b
\]
on that subsequence \cite[A charged disk throughout the subcritical
band]{Baker2026ChiralMaster}. For every sufficiently late charged
index \(d_i<a\) and \(e_i<2T\). Corollary~\ref{cor:fixed} then
rules out a round umbilic face at each of those actual disks.
The constants \(C_i\) need not be uniform. The budget compares
positive metrics at positive widths at each fixed index before
using the limiting bound.

\begin{corollary}[Defect on precompact charged collars]\label{cor:charged}
On any subsequence of those original marked charged collars having
the precompactness hypothesis of Theorem~\ref{thm:compact},
the face defects have strictly positive lower limit, provided the
required subwidths fit in the permitted regions.
\end{corollary}

\begin{proof}
If the lower limit were zero, pass to a further subsequence with
defect tending to zero and apply Theorem~\ref{thm:compact}. Along
the resulting paths the joined peaks are at most \(e_i+o(1)<2T\)
for every sufficiently late index. Their marked capped exteriors
are positive, and in the selected sign their energies on \(M\)
are \(e_i-m_i^{\sigma_*}+o(1)\le e_i-a+o(1)\). At a sufficiently
late finite index this is strictly below \(I\), contradicting
\eqref{eq:disk-budget}. The argument gives no numerical lower
bound on the defect without further information on the family.
\end{proof}

\subsection{A separate test on an enlarged disk}

The master also constructs, for each fixed bisection count
\(k\ge1\), an enlarged parametrized disk \(B_i\) protecting an
outer disk \(Q_i\) that contains the original marked bicollar.
Write \(q_i^\sigma=E_\sigma(g_i;Q_i)\),
\(v_i^{B,\sigma}=E_\sigma(g_i;M\setminus\operatorname{int}B_i)\),
and \(b_i^\sigma=e_i-v_i^{B,\sigma}\). Its four-measure
selection gives, for fixed \(k\),
\[
 v_i^{B,\sigma}=2^{-k}(e_i-q_i^\sigma)+r_i^\sigma,
 \qquad r_i^\sigma\longrightarrow0
\]
in both signs \cite[A necessary disk-side saving]{Baker2026ChiralMaster}.
This is a statement about the raw measures and the actual enlarged
marking. If a proposed path ends in a round product germ there
and has strict joined peak \(H_i<2T\), the exterior cap still
has energy at least \(I\). The signed exterior increment is
therefore at least \(b_i^\sigma-d_i\), and the saving
\(S_i^\sigma=b_i^\sigma-E_\sigma(h_i;B_i)\) on the entire
disk satisfies
\[
 S_i^\sigma\ge b_i^\sigma+I-H_i>b_i^\sigma+I-2T .
\]
For \(k=1\) and \(I=4T/3\), the last bound has lower limit at
least \(a/2\) in sign \(\sigma_*\), using \(q_i^{\sigma_*}\ge a\).
More generally the limiting lower bound in either sign is at
least \((2-2^{-k})I-2T\). If the master's non-effective gap on
the given nonstandard type is \(I>T+\varepsilon_*\), a fixed
\(k\) with \(2^{-k}(T+\varepsilon_*)\le\varepsilon_*\) gives a
limiting lower bound greater than \(\varepsilon_*\)
\cite[Uniform recognition by the Weyl spectrum]{Baker2026ChiralMaster}.

A change confined to the final disk-side seam saves at most
the incoming chiral mass of that seam. Thus a sequence of
seams with vanishing raw disk-side mass cannot furnish the
required savings in the positive lower bound regimes just
listed. If the whole \(Q_i\) is preserved up to conformal
scale, then \(H_i\ge I+\max_\sigma q_i^\sigma\), since the
enlarged disk retains at least that core energy in its
respective sign and the exterior cap consumes at least \(I\).
These are distinct finite-source conditions on a matched
comparison.

\paragraph{An interior change on a disjoint support.}
Let a compactly supported interior preparation on the original disk
follow a positive path of nonincreasing total energy from \(e_i\) to
\(e_i-S\). Signature gives the same saving on either chirality
and \(0\le S\le d_i\). Let a specified seam path have raw joined peak
\(H_C\) and meet the support of the preparation in no point,
including the transitions at the native widths. Locality gives the
peak after performing them in this order as
\(\max(e_i,H_C-S)\). The complete capped exterior is exactly the
one produced by the seam path. If the combined peak is below \(2T\)
and \(c^\sigma\) is the raw disk-side mass of the seam support, the
whole disk-side saving at the matched endpoint is at most
\(d_i+c^\sigma\). This cannot pay the balanced lower bound above
when \(c^\sigma\to0\) at sufficiently large fixed depth. For
simultaneous disjoint deformations, suppress each support in turn.
At the first isolated crossing of \(2T\) the other isolated energy
is at least \(I\) by continuity from the positive prefix. The
combined peak is then at least \(2T-d_i\). If neither isolated
path crosses, the same disk-saving obstruction applies. Hence a
matched, strictly subcritical simultaneous path on that balanced
family would have to enter \([2T-d_i,2T)\). All these estimates
use the actual complete support of each change.

\section{A curved Euclidean filling on the original disk}\label{sec:convex}

The replacement of a parametrized disk can use a curved boundary of a
flat four-ball. The angular metric and its trace-free normal derivative
determine whether a strictly convex Euclidean ball supplies the required
first jet. The scalar curvature of the retained exterior enters a
separate test for positivity.

\begin{theorem}\label{thm:convex}
Let \(g\) be a smooth metric on the oriented smooth homotopy sphere
\(M\), let \(D\) be a smoothly parametrized closed disk, and put
\(V=M\setminus\operatorname{int}D\).
On its actual boundary \((\Sigma,h)\) write \(K\) for the shape
endomorphism with normal pointing from \(D\) into \(V\), and
\(H=\tr K\). Define
\[
 G=\tfrac12 R_h\operatorname{Id}-\Ric_h^\sharp,\qquad
 S=(\det G)^{1/2}G^{-1}\quad\text{when }G>0 .
\]
Suppose \(G>0\), \((\nabla_X S)Y=(\nabla_Y S)X\), and
\(S^\circ=K^\circ\). Then
\(E_+(g;D)=E_-(g;D)\). If, in addition, \(R_g>0\) on \(V\)
and \(H\le\tr S\), then
\[
 I^+(M)\le E_\sigma(g;V),\qquad
 E_\sigma(g;D)\le E_\sigma(g;M)-I^+(M)
 \quad(\sigma=+,-).
\]
\end{theorem}
\begin{proof}
In an orthonormal eigenbasis, the eigenvalue of \(G\) at the first
axis is the sectional curvature of the complementary plane, and
cyclically. The eigenvalues of \(S\) thus have pairwise products
equal to those sectional curvatures. In dimension three these relations
give the whole Gauss equation since the curvature operator is diagonal
in the corresponding planes. Codazzi and simple connectivity give
a Euclidean immersion \(F\) with metric \(h\) and shape \(S\)
in our outward convention \cite[Theorem 1.3]{LiSimonZhaoHu2015}.
Its unit normal map is a covering of the unit three-sphere and hence
a diffeomorphism. Each height in a unit direction has precisely two
critical points and its unique maximum has that outward normal, by
positivity of \(S\). Each supporting plane of the convex hull
therefore contacts a single image point, and each image point has
a unique preimage by this maximum test. The hull has nonempty
interior and boundary \(F(\Sigma)\). Radial projection from an
interior point identifies its smooth strictly convex domain with a
standard ball (interpolate radially from a small concentric ball
and use polar coordinates only away from the center).

Let \(b\) be its flat metric. Choose smooth \(u>0\) on the ball
with \(u=1\) at the boundary and
\(\partial_{\nu_b}u=\alpha=(H-\tr S)/3\) in the outward normal.
For example prescribe these jets to \(\log u\) in a collar, then
extend by a cutoff. The cap \(u^2b\) has shape \(S+\alpha
\operatorname{Id}=K\) there. Glue by \(F\) on the marked face,
using its inverse in the other direction. The angular discrepancy
between this ball and the original disk chart extends over a
standard ball by the sphere diffeomorphism theorem
\cite{Hatcher1983Smale}. This identifies the assembly with the
original smooth \(M\).

Use signed Gaussian distance \(t\) in each metric, increasing into
\(V\), with the same boundary parameter. At zero both tangential
families have value \(h\) and derivative \(2h(K\cdot,\cdot)\).
The original physical metric supplies a smooth extension of the
exterior family to negative \(t\). In a cap strip of width \(w\)
interpolate to that extension with a fixed smooth cutoff, equal
to the extension near zero and to the cap near \(-w\). The difference
of the two families is \(O_g(t^2)\), its first derivative is
\(O_g(|t|)\), and its second derivative is bounded at this fixed
source and fill. The same Taylor bounds hold for tangential derivatives
as needed. The \(w^{-1}\) and \(w^{-2}\) cutoff terms therefore give
bounded curvature. Take \(w\) inside both the original inward
Gaussian radius and the cap radius. The exterior is unchanged on its
entire native side, and the cap outside this connector is conformally
flat. For the resulting smooth metric \(p_w\) on \(M\),
\[
 E_\sigma(p_w;M)=E_\sigma(g;V)+O_{g,b,u}(w).
\]
Signature on \(p_w\) and on \(g\) gives the claimed disk mass equality
on sending \(w\downarrow0\).

For the sign assertion consider the limiting joined quadratic form.
Integration by parts on the cap, using
\(R_{u^2b}=-6u^{-3}\Delta_b u\), gives for a common trace \(f\)
\[
 Q(f)=\int_V(6|df|_g^2+R_g f^2)\,dV_g+
 6\int_D |d(uf)|_b^2\,dV_b+
 2\int_\Sigma(\tr S-H)f^2\,dA_h .
\]
Under the assumed signs this form controls the \(H^1\) norm on
the fixed joined space. Indeed the exterior term controls the
exterior norm and hence its boundary trace. On the ball the gradient
of \(uf\) and its boundary trace control its full \(H^1\) norm by
Poincar\'e's inequality with boundary term. Multiplication by the
fixed positive smooth \(u\) is an isomorphism there. The metrics
\(p_w\) converge uniformly to the piecewise metric. Their scalar
potentials differ from its a.e. scalar potential only on a strip
of volume \(O(w)\) and remain uniformly bounded there. H\"older's
inequality and the fixed compact four-dimensional Sobolev inequality
give a form error \(o(1)\|f\|_{H^1}^2\), including the change of
metric density and inverse. All sufficiently small positive widths
therefore give positive classes. Taking their infimum proves the result.
\end{proof}

This argument also works when the displayed boundary form is coercive
by direct measurement. It proves positivity at the new endpoint.
The joined peak of a continuous deformation replacing the entire disk
has not been estimated. The ball is replaced in full, so a requirement
to preserve an interior core demands a different comparison. At a
sufficiently late original charged index the conclusions exclude the
joint boundary hypotheses in the positive Yamabe representative.
The shape tests require the actual induced metric and second fundamental
form at that index.

\paragraph{An induced height slice.}
Let \(z\) be a regular height of a Euclidean hypersurface immersion
in \(\mathbb R^5\). Suppose the slice in the height hyperplane
is an embedded three-sphere of positive definite outward shape, and use its own flat
interior as the candidate ball. Along the level the immersed unit
conormal of the four-manifold has the form \(a v+b n\), where
\(v\) is the constant height vector, \(n\) the outer Euclidean
normal of the slice in that hyperplane, \(a^2+b^2=1\) and
\(|a|=|\nabla z|>0\). Differentiation against any slice tangent
gives \(K=bS\) for the induced four-metric. Hence \(H<\tr S\), but
\[
 |K^\circ-S^\circ|=(1-b)|S^\circ|
 \ge \tfrac12 |\nabla z|^2 |S^\circ|.
\]
Exact matching throughout this convex slice forces \(S^\circ=0\).
Codazzi then makes \(S\) a positive constant multiple of the identity,
so the slice is round. A transverse immersed height therefore gives a
specific first-jet obstruction on its own nonround convex sections.
An intrinsic deformation or a different comparison gauge would
require new boundary measurements.

\section{Finite errors at a curved face}\label{sec:finite-face}

Keep the convex Euclidean filling and the actual marking of
Theorem~\ref{thm:convex}, now allowing
\(D_0=h((K^\circ-S^\circ)\cdot,\cdot)\ne0\).
Prescribe the same conformal normal derivative on the cap so that
the two mean traces agree. Let \(a_t,p_t\) be the tangential
families in their paired Gaussian coordinates, with the original
exterior extended into the disk by its own metric. They agree in
value and \(a'_0-p'_0=2D_0\). For \(-w\le t\le0\) set \(x=-t/w\)
and interpolate by \(q_t=p_t+\chi(x)(a_t-p_t)\), with \(\chi=1\)
near zero and zero near one. Take \(w\) within both actual inward
radii. The complete exterior remains in its original metric.

\begin{proposition}\label{prop:finite-curved}
Let \(\psi=x\chi\) and write \(c_\sigma(w)\) for the entire new
cap energy, supported in this one strip. At each fixed source and
fill there are bounded curvature remainder constants \(B_\sigma(w)\)
for sufficiently small \(w>0\) such that, with
\[
 L=\frac{\|\psi''\|_{L^2(0,1)}\|D_0\|_{L^2(h)}}{2\sqrt w},
 \qquad b_\sigma=B_\sigma(w)\sqrt{w\operatorname{area}_h(\Sigma)},
\]
one has \((L-b_\sigma)_+^2\le c_\sigma(w)\le(L+b_\sigma)^2\).
The scalar curvature stays bounded on this strip. If the limiting
boundary form of Theorem~\ref{thm:convex} is coercive, all sufficiently
small positive widths give positive capped exterior metrics and
\[
   I\le e-m^\sigma+c_\sigma(w),\qquad
   c_+(w)-c_-(w)=m^+-m^- .
\]
\end{proposition}
\begin{proof}
The sole second derivative term growing as \(w^{-1}\) is
\(q''_t=-2\psi''D_0/w+O(1)\). In Gaussian coordinates its normal
curvature block is \(\psi''D_0/w+O(1)\); all tangential and mixed
curvature blocks are bounded. In each chiral Weyl block this
trace-free normal block enters with coefficient \(1/2\). The
change to orthonormal frames and the square root of the volume
density relative to \(dt\,dA_h\) differ from their face values by
\(O(w)\). Bound the residual matrix in pointwise Hilbert--Schmidt
norm by \(B_\sigma(w)\) and apply the two triangle inequalities
in \(L^2\). One can compute these remainders from the two native
tangential families by subtracting the displayed matrix from the
density-weighted Weyl blocks. Taylor bounds on their second jets
over a buffered interval also bound them from above.

In the scalar contraction \(\tr_{q_t}D_0=O(w)\).
The scalar potential error relative to the piecewise form has
bounded size on a strip of width \(w\); its integral against
\(f^2\) is \(O(w)\|f\|_{H^1}^2\) by the uniform slice trace
inequality in a fixed collar. The principal part changes by
\(O(w)\) in quadratic form norm. The boundary form is exactly
the one from integration by parts on the two pieces, since the
mean traces agree. Coercivity persists at positive widths.
The two final statements follow by complete-side accounting,
the fixed-type infimum and closed signature.
\end{proof}

Thus a nonzero fixed error costs
\(\|\psi''\|_2^2\|D_0\|_2^2/(4w)+O(1)\) in \emph{each} sign.
For a finite test one needs a width in the admissible chart for
which coercivity survives and the upper cost is below \(m^\sigma-d\)
in a sign. The signed difference alone cannot price this comparison.

\paragraph{What the signed difference measures.}
Write \(\delta(g;D)=E_+(g;D)-E_-(g;D)\). Let a reference ball
\((B,b)\) induce exactly \(h\) via an oriented boundary marking, and
denote the two outward shapes by \(K,S\). View them as
\(T\Sigma\)-valued one-forms \(k,s\), set \(a=(k-s)^\circ\) and write
\(\mathsf D=d_{\nabla^h}\). Then
\[
 \delta(g;D)-\delta(b;B)
 =\int_{\partial D}\bigl(\langle a\wedge\mathsf D a\rangle+
                 2\langle a\wedge\mathsf D s\rangle\bigr).
\]
Indeed the signed density is \(-\tr(F\wedge F)/2=4\pi^2p_1\).
The connection at the face differs from the product connection
by normal-tangent off-diagonal entries \(k,-k^{\mathsf T}\).
Transgression gives \(\langle k\wedge\mathsf D k\rangle\) in
outward-first orientation. Apply the same calculation to the reference
and glue product germs. The topological term is
\(12\pi^2\tau(D\cup(-B))=0\) for balls. Symmetry and torsion-freeness
remove the scalar trace, and integration by parts on the closed
face gives the formula
\cite[\S3]{ChernSimons1974Characteristic}\cite{Hirzebruch1971Signature}.
Codazzi kills the mixed term for a space form. On a Euclidean
hypersurface height with \(K=\beta S\), the flux against its own
Euclidean fill vanishes even when the two trace-free shapes differ.
In fact the Gauss curvature operator of a Euclidean hypersurface is
diagonal on principal two-planes, so the complementary chiral
blocks have the same Weyl spectrum pointwise.

The cutoff coefficient has a sharp lower envelope within this
fixed-profile construction. The conditions on \(\psi\) are
\(\psi(0)=\psi(1)=\psi'(1)=0\), \(\psi'(0)=1\).
The cubic \(x(1-x)^2\) minimizes \(\int|\psi''|^2\) with
these four data: integration by parts annihilates the cross term
with every clamped variation. Its integral is four. Replacing the
cubic by \(x\) near zero and zero near one on shrinking end
intervals changes its \(H^2\) norm by \(o(1)\), by its quadratic
contact at both ends. Thus even optimizing fixed clamped profiles
leaves the leading coefficient at least \(\|D_0\|_2^2/w\).
This is a bound for the specified interpolation and does not
optimize over general metrics.

For a flat convex reference write
\(\Delta=\delta(g;D)\). In the flux integral one can replace
\(\mathsf D a\) by \(\mathsf D k\), since \(\mathsf D s=0\)
and terms from the trace pair to zero. In particular
\[
 |\Delta|\le \|D_0\|_2\|\mathsf D k\|_2 .
\]
A measured imbalance and finite nonzero Codazzi norm now give a
lower bound on \(L\) in Proposition~\ref{prop:finite-curved}
by substitution of \(|\Delta|/\|\mathsf D k\|_2\) for
\(\|D_0\|_2\). In a physical Gaussian strip \(U\) of width \(w_0\)
the Codazzi curvature occupies one of two transposed mixed blocks.
The ordered tensor norm has squared norm four times the squared
curvature-operator norm. Coarea therefore gives a parallel face with
\(\|\mathsf D k\|_2^2\le \int_U|\Rm_g|_{\rm tens}^2/(8w_0)\).
The boundary isometry, imbalance and common inward interpolation
width must belong to this very face for these conclusions to combine.

\subsection{A hyperbolic reference for an induced height}
A change in reference curvature can repair the trace-free mismatch
on a specific nonround height. For a strictly convex Euclidean
section of shape \(S\), a fit at sectional curvature \(\kappa\)
has candidate shape \(\beta S+c\operatorname{Id}\) where
\(K=\beta S\), \(|\beta|<1\). Subtracting the Gauss equations for
pairs of eigenvalues \(s_i\) yields
\((s_j-s_l)((1-\beta^2)s_i-\beta c)=0\).
Three distinct principal curvatures preclude a fit. If the repeated
one is \(u\), a positive nonumbilic fit requires \(\beta>0\) and
\[
 \kappa=-k^2=-\frac{(1-\beta^2)u^2}{\beta^2},\qquad
 c=\frac{k^2}{\sqrt{u^2+k^2}}.
\]
Here \(k\) must be constant on the face and the reconstructed shape
must be positive. For the ellipsoid
\[
 |x|^2/A^2+y^2/B^2+z^2/C^2=1\quad(x\in\mathbb R^3),
 \qquad D=\{z\ge z_0\},\quad 0<z_0<C,
\]
the repeated principal curvature of the section is
\(u=1/(A^2|(x/A^2,y/B^2)|)\) and
\(\beta=u/\sqrt{u^2+k^2}\), \(k=C^2/(A^2z_0)\).
For completeness write the Euclidean meridian as \((r(s),y(s))\)
in arclength, with \(u=y'/r\). In axial hyperbolic coordinates the
warped coefficients of \(dZ^2,g_{S^2}\) are respectively
\(\cosh^2(k\rho),k^{-2}\sinh^2(k\rho)\). Set
\[
 \rho=k^{-1}\operatorname{arsinh}(kr),\qquad
 Z'=\frac{r\sqrt{u^2+k^2}}{1+k^2r^2}.
\]
The induced metric is \(ds^2+r^2g_{S^2}\); the poles are smooth by
the parity of the ellipsoid meridian. The axial coordinate is strictly
increasing in the interior. The angular principal curvature in the
outward normal is \(\sqrt{u^2+k^2}\), and the meridional curvature
follows from Gauss, giving precisely \(\beta S+c\operatorname{Id}\).
This explicit embedded profile bounds a ball.

Write the hyperbolic metric in centered Poincar\'e coordinates as
\(\Omega^2g_{\rm E}\), and normalize the cap by \(\mu^2\) with
\(\mu=1,\ \partial_\nu\log\mu=-c\) on the face. The boundary
term remaining after integration by parts on the cap is
\(6\int_\Sigma(c-\partial_\nu\log\Omega)f^2\), in addition to
\(6\int |d(\mu\Omega f)|^2_{g_{\rm E}}dV_{g_{\rm E}}\).
For a round face of intrinsic radius \(R\) the difference in parentheses
is \(k^2R/[\sqrt{1+k^2R^2}(1+\sqrt{1+k^2R^2})]>0\).
The inequality persists for nearby ellipsoids. The entire exterior
has positive scalar curvature. Thus the one-sided exact-jet smoothing
gives a positive capped exterior in the actual induced gauge.

Keep \(z_0/C\) fixed and take \(B,C\) close to \(A\), each different
from \(A\). On \(y=0\) the four principal curvatures are
\((p,p,q,v)\), with \(p\ne q,p\ne v\) away from the axes; the
complementary Weyl entry is \((p-q)(p-v)/3\).
Both complete sides have positive mass in both signs.
At the round parameter the intrinsic disk and the normalized cap
can be made identical by radial conformal coordinates. Transport this
choice smoothly and correct the normal jet in a fixed collar. With
\(\epsilon=|B/A-1|+|C/A-1|\), the paired tangential difference is
\(O(\epsilon t^2)\) through its differentiated Taylor bounds.
Clamping inside a common native width leaves a uniform
\(O(\epsilon)\) second-jet bound relative to round. Affine
interpolation of the two metrics on the fixed marked sphere keeps
positive scalar curvature for small \(\epsilon\). Its whole peak in
either sign is at most \(E_\sigma(g;V)+O(\epsilon^2)\);
the endpoint adds only \(O(\epsilon^2w)\) to that exterior.
At a fixed charged member one positive width therefore saves energy
in both signs. The entire disk is replaced. These are induced
metrics on the standard sphere. The selected minimizing disk has
its own marking; no supplied calculation identifies its shape or
its boundary coercivity with this one.

\section{Testing a coupled path on the original marking}\label{sec:isolate}
Let \(g\) be the incoming positive source and \(h_t\) a proposed
matched path on that same sphere with \(h_0=g\) and full joined peak
\(H<2T\). Write \(e_t=E_+(h_t;M)\), \(d=e-I\). Choose an actual
disk-side collar \(A\) up to the face and a smooth cutoff \(\chi\)
from zero at its inner edge to one at the face, constant near both.
Extend it by zero on the deep disk and by one on all of \(V\).
Set \(p_t=g+(1-\chi)(h_t-g)\) and \(q_t=g+\chi(h_t-g)\).
These are closed smooth metrics on the same marking. Write
\[
 \begin{split}
 J^\sigma(t)&=E_\sigma(p_t;A)+E_\sigma(q_t;A)
                    -E_\sigma(h_t;A)-E_\sigma(g;A),\\
 R^\sigma(t)&=E_\sigma(p_t;A)-E_\sigma(h_t;A).
 \end{split}
\]
The complete closed energies of the two isolated paths sum to
\(e+e_t+J^\sigma(t)\), so \(J^+=J^-\). Also
\(E_\sigma(p_t;M)=e-m^\sigma+E_\sigma(h_t;D)+R^\sigma(t)\).
At the matched round-product endpoint Lemma~\ref{lem:cap} supplies
a positive Weyl-free cap in the original disk marking, so the complete
exterior forces
\(E_\sigma(h_1;V)\ge I\), including all exterior connectors.
If \(p_1\) is positive this gives
\[
 R^\sigma(1)\ge m^\sigma+I-e_1-d\quad(\sigma=+,-).
\]
Otherwise one of the two isolated paths reaches \(2T\). At their
first such time \(s\), the other still has energy at least \(I\)
by the scalar-flat wall and continuity. Reaching the energy level
need not make the isolated metric scalar-flat. Hence
\[
 J^+(s)=J^-(s)\ge2T-d-e_s\ge2T-d-H.
\]
These alternatives apply also when both changes use the same
physical strip. For an enlarged disk substitute exactly
\(m^\sigma=(1-2^{-k})e+2^{-k}q^\sigma-r^\sigma\).
If the protected core is kept up to conformal scale, also
\(e_1\ge I+\max_\sigma q^\sigma\).
At large fixed depth on a late sequence with \(I>T\), the
endpoint bound is strictly positive. For \(H\le2T-\eta\)
the crossing bound costs at least \(\eta-d\).
A raw small incoming collar does not bound either isolated
connector. For example a native cutoff derivative of order
\(w^{-2}\) multiplies the full tensor \(h_t-g\) in its curvature.
The isolation test requires measured metrics on the whole path.


\section{Outstanding challenges}

The signed second-jet coefficient and the mode estimates price a
particular intrinsic interpolation. They leave open optimization of
the path, the cutoff and the physical coordinates. In particular,
a small integrated face defect by itself leaves available widths,
ellipticity and normal jet bounds undetermined. The compact result
uses a positive definite smooth limit across an actual common marked
slab. Establishing an applicable compactness or jet estimate on the
selected physical collars requires control of their scale, chart
distortion, protected regions and normal geometry.

The marked handle criteria concern specified cuts with their full
gluing data \cite{Baker2026HandleCuts}. Using their geometry in
the charged comparison requires a common physical realization or a
controlled connecting region on the same metric. Likewise a repair
of an enlarged charged disk requires a decrease of energy on the
whole disk sufficient for the finite inequalities above. The
arguments give conditional exclusions on entered positive subcritical
types. Producing such entry or a recognizing metric on an arbitrary
smooth homotopy four-sphere is an additional global problem
\cite{Baker2026CollarProgram,Baker2026ChiralMaster}.

\bibliographystyle{unsrturl}
\bibliography{references}
\end{document}
